\section{基础设施与 DevOps 原则}

\subsection{监控的四个黄金信号}

\begin{importantbox}{Google SRE 的四个黄金信号}
\begin{enumerate}
    \item \textbf{Latency}（延迟）：请求的响应时间
    \item \textbf{Errors}（错误）：失败请求的比率
    \item \textbf{Traffic}（流量）：每秒请求数（requests/sec）
    \item \textbf{Saturation}（饱和度）：资源利用率，接近极限时系统性能急剧下降
\end{enumerate}
此外，还需要监控生产环境的 traces（请求追踪链路）。
\end{importantbox}

\subsection{事件响应实战：凌晨 3:12 的数据库 500 错误}

当 PagerDuty 告警显示数据库查询出现 500 错误激增时，标准的响应流程（playbook）如下：

\begin{enumerate}
    \item \textbf{确认与评估}（Acknowledge and assess）
    \item \textbf{检查数据库 + 应用}（Check DB + app）
    \item \textbf{识别近期变更}（Identify recent changes）
    \item \textbf{定位影响范围}（Localize blast radius）
    \item \textbf{执行缓解措施}（Execute mitigations）
    \item \textbf{稳定 + 监控}（Stabilize + monitor）
    \item \textbf{沟通}（Communicate）
    \item \textbf{文档记录}（Document）
\end{enumerate}

\subsection{关键度量指标}

\begin{itemize}
    \item \textbf{MTTR}（Mean Time to Repair）：平均修复时间
    \item \textbf{参与事件的工程师数量}
    \item \textbf{客户 SLA 合规报告}
\end{itemize}

\subsection{本章小结}

四个黄金信号（延迟、错误、流量、饱和度）是监控生产系统的基础框架。事件响应遵循标准化的 playbook，关键度量是 MTTR 和影响范围。

